% Generated draft: requires mathematical and source review.
\documentclass[11pt]{article}
\usepackage[margin=25mm]{geometry}
\usepackage{fontspec}
\setmainfont{Latin Modern Roman}
\usepackage{amsmath,amssymb,mathtools,unicode-math}
\setmathfont{Latin Modern Math}
\usepackage{enumitem,needspace,xurl,hyperref}
\setlength{\parindent}{0pt}
\setlength{\parskip}{5pt}
\setlength{\emergencystretch}{3em}
\newcommand{\N}{\mathbb{N}}
\newcommand{\Z}{\mathbb{Z}}
\newcommand{\Q}{\mathbb{Q}}
\newcommand{\R}{\mathbb{R}}
\newcommand{\C}{\mathbb{C}}
\newcommand{\F}{\mathbb{F}}
\newcommand{\A}{\mathbb{A}}
\newcommand{\PP}{\mathbb{P}}
\newcommand{\E}{\mathbb{E}}
\newcommand{\eps}{\varepsilon}
\newcommand{\dd}{\,\mathrm{d}}
\begin{document}
\section*{UnsolvedMath ID: 30002811}
\textbf{Source problem code:} catalogue entry
\subsection{Definitions and mathematical statement}
The supplied problem concerns a discrete-time control system
\[
 x(k+1)=f(x(k),u(k))
\]
with stage cost \(\ell(x,u)\). The state and control domains, admissibility constraints, optimization horizon, terminal conditions, and precise meaning of the turnpike property are not specified in the record.

To express the supplied dissipativity condition, write \(X\) for the state domain in an unspecified normed space, \(U\) for the control domain, and \(D\subseteq X\times U\) for the admissible state--control pairs, with \(f:D\to X\) and \(\ell:D\to\R\). These symbols describe data that must be supplied; they do not impose a particular state space or admissibility model. An admissible equilibrium is a pair \((x_e,u_e)\in D\) satisfying \(f(x_e,u_e)=x_e\). The class \(\mathcal K_\infty\) consists of continuous, strictly increasing functions \(\rho:[0,\infty)\to[0,\infty)\) such that \(\rho(0)=0\) and \(\lim_{r\to\infty}\rho(r)=\infty\).

The source's strict dissipativity condition asks for an equilibrium \((x_e,u_e)\), a storage function \(\lambda:X\to\R\), and \(\rho\in\mathcal K_\infty\) such that
\[
 \ell(x,u)-\ell(x_e,u_e)
 +\lambda(x)-\lambda\bigl(f(x,u)\bigr)
 \geq \rho\bigl(\lVert x-x_e\rVert\bigr)
 \qquad\text{for every }(x,u)\in D.
\]
Here universal quantification over admissible pairs makes explicit the usual reading of the displayed source condition; the intended admissible set and whether a local restriction is intended remain unspecified. The record gives no regularity or boundedness requirements on \(\lambda\), and the strictness term measures state deviation only.

\textbf{Problem.} Determine assumptions under which a turnpike property for the stated optimal-control problem implies strict dissipativity in the displayed sense. A precise converse theorem requires specifying the turnpike property, its quantifiers and uniformity, the optimization problem and admissible initial states, and its relationship to the equilibrium \((x_e,u_e)\). These choices cannot be recovered uniquely from the supplied record. In particular, the source does not explicitly require \((x_e,u_e)\) to minimize the stage cost among admissible equilibria.

The record describes strict dissipativity implying a turnpike property as the known direction, without supplying its hypotheses or identifying the turnpike variant. This is retained as a qualified source assertion, not an unconditional theorem.
\subsection{Short English statement}
The problem asks when optimal trajectories spending most of their horizon near an equilibrium force a strict dissipativity inequality relative to that equilibrium. This description of turnpike behavior is informal: the source does not specify its exact definition or the control problem's constraints and boundary conditions. The supplied status is \emph{partially solved}; a precise theorem statement requires resolving these omissions.
\subsection{Sources}
\begin{itemize}
\item Supplied local UnsolvedMath registry record, permanent ID \(30002811\), source code \texttt{catalogue entry}, title \emph{Turnpike Properties as Criteria for Strict Dissipativity}, in the collection \emph{Oberwolfach Reports Open Problems}. Project registry references: \texttt{data/unsolvedmath/problems.json} and supplementary \texttt{catalog/problems.json}; neither file was independently inspected for this draft.
\item \emph{Control Theory: A Mathematical Perspective on Cyber-Physical Systems} (2015), Oberwolfach Reports, pp.~597--659, DOI \texttt{10.4171/owr/2015/12}. Supplied report reference and literature-source link: \url{https://doi.org/10.4171/owr/2015/12}. The record attributes the problem to workshop open-problem discussions and supplies the proposed year 2015, but names no individual proposer and gives no problem-specific page.
\item Supplied literature reference labelled ``Grune and Muller, strict dissipativity'': \url{https://doi.org/10.1016/j.sysconle.2016.01.003}. The record supplies no fuller bibliographic citation.
\item Literature assessment dated \textbf{2026-08-21}: the supplied curation states that, under controllability and structural assumptions, strict dissipativity and several turnpike variants are equivalent, while a converse theorem under minimal general assumptions remains unavailable. This is supplied dated curation, not an independently verified literature assessment. Its time-sensitive status and cited claims require independent verification.
\item The registry labels the problem \texttt{partially\_solved} and the statement \texttt{corrected\_verified}, with the explanation that mathematical content was preserved while notation, typography, or omitted definitions were repaired for publication. These are supplied registry labels, not mathematical verification performed for this draft.
\end{itemize}
\end{document}
